\section{A curve with prescribed residue data}
\label{curve:section}

We use additive notation for line bundles.  If $B$ is a line bundle and
$D$ is a divisor, $B+D$ means $B\otimes\OO(D)$; restrictions of bundles
will occasionally be left implicit.  We first construct a marked curve
and a fourth root of its canonical bundle.  An auxiliary elliptic curve
will control the directions and the normalization of the residue data.
The output is collected in Proposition~\ref{prop:curve-data}: besides
$L^4\simeq K_C$, we need precise section dimensions and evaluation maps
at six marked points. These conditions will make the low-order gluing
obstructions computable.

\subsection{The elliptic model and its quadratic form}

Let $\Gamma$ initially be a smooth elliptic curve.  Choose a degree-two
line bundle $F$ and a line bundle $D_0$ of exact order four, and set
$S=F+D_0$.  Both $F$ and $S$ are generated by their two-dimensional
spaces of sections.  The product of their associated maps embeds
$\Gamma$ in
\[
  \Sigma=\PP^1\times\PP^1.
\]
Indeed, after choosing an origin on $\Gamma$, the two fiber involutions
have the form $p\mapsto a-p$ and $p\mapsto a+d_0-p$, where $d_0\ne0$
represents $D_0$.  Two distinct points cannot be interchanged by both
involutions, and a point cannot be fixed by both.  The product map is
therefore injective and immersive, hence is a closed embedding.  Its
two projection degrees show that its image has bidegree $(2,2)$.
We now also write $F=(1,0)$ and $S=(0,1)$ for the ruling bundles on
$\Sigma$.

On $\Gamma$ put
\[
  \Delta=2D_0,\qquad \Phi=F+\Delta,\qquad M=3\Phi.
\]
Thus $\Delta$ is nontrivial of order two, and
\[
  \Gamma|_\Gamma=2F+2S=4F+\Delta,
  \qquad M=3F+\Delta,\qquad 2M=6F.
\]
Consequently
\begin{equation}\label{curve:bundle-identities}
\begin{split}
  4M&=(2\Gamma+4F)|_\Gamma,\\
  2M-(\Gamma+F)|_\Gamma&=\Phi,\\
  (\Gamma+2F)|_\Gamma-2M&=\Delta.
\end{split}
\end{equation}
Fix isomorphisms in these identities compatibly.  In particular,
viewing $\Phi$ as $2M-(\Gamma+F)|_\Gamma$, the first identity induces
the identification $2\Phi=2F|_\Gamma$ that we use below.

Since $K_\Sigma+\Gamma=0$, choose a nowhere vanishing logarithmic
two-form $\Omega$ with pole divisor $\Gamma$.  Its residue
$\omega=\Res_\Gamma\Omega$ is a nowhere vanishing differential on
$\Gamma$.  Our residue convention places the normal differential
first: locally $\Res_{y=0}((dy/y)\wedge dz)=dz$.
Choose a basis $u=(u_1,u_2)$ of $H^0(\Sigma,F)$, and write
\[
  \eta_\Gamma=u_1\,du_2-u_2\,du_1
       \in H^0(\Gamma,K_\Gamma+2F).
\]
The expression is independent of the local frame for $F$: in that
frame it is the indicated coefficient differential times the square
of the frame.

There is a unique element of $\Sym^2H^0(\Gamma,\Phi)$ whose image under
multiplication is
\begin{equation}\label{curve:quadratic}
  -\frac{\eta_\Gamma}{\omega}\in H^0(\Gamma,2\Phi)
       =H^0(\Gamma,2F).
\end{equation}
It defines a nondegenerate quadratic form $Q$ on
$H^0(\Gamma,\Phi)^*$.  To see both existence and nondegeneracy, let
$\iota_F$ and $\iota_\Phi$ be the two degree-two involutions.  One has
$\iota_\Phi=t_\Delta\iota_F$, where $t_\Delta$ is translation by the
nonzero point of order two corresponding to $\Delta$.  The four
ramification points of $F|_\Gamma$, namely the fixed points of
$\iota_F$, are paired without fixed points by $\iota_\Phi$.  They form
two distinct reduced fibers of the $\Phi$-map.  Their divisor is
therefore the pullback of the zero divisor of a binary quadratic with
two distinct roots.  Scaling this quadratic gives exactly
\eqref{curve:quadratic}.  Uniqueness follows because pullback from the
surjective $\Phi$-map is injective on sections of $\OO_{\PP^1}(2)$.

For $q\in\Gamma$, evaluation in a local frame of $\Phi$ gives a
nonzero element of $H^0(\Gamma,\Phi)^*$, defined up to a nonzero scalar.
The corresponding projective direction will be denoted by
$[\operatorname{ev}_q]$.  By \eqref{curve:quadratic}, this direction is
isotropic precisely when $q$ is a ramification point of $F|_\Gamma$.
For any quadratic form in this section, we use its polar form
\[
 x\mathbin{\cdot}y=\frac{Q(x+y)-Q(x)-Q(y)}2,
 \qquad x^2=Q(x).
\]

The six vectors below encode both a pointwise normalization and a total
residue.  If $w=\sum_jw_j$ and $w_j^2+w_j\mathbin{\cdot}w=1$ at
each marked point, then
\[
 \sum_{j=1}^6(w_j^2-1)=-w^2.
\]
The curve we construct has $\deg L=10$.  Thus the choice $w^2=90$
will give the total $-9\deg L$ required by the order-two obstruction
in Section~\ref{jets:section}, while retaining all six pointwise
normalizations.

\begin{lemma}\label{curve:six-vectors}
In a two-dimensional nondegenerate complex quadratic space, one can
choose three distinct nonisotropic directions, avoiding any prescribed
finite set of projective directions, and nonzero vectors
$w_1,\ldots,w_6$ that use each of these directions twice, with the two
vectors in each pair equal, such that
\begin{equation}\label{curve:vector-conditions}
  w=\sum_{j=1}^6w_j,\qquad
  w^2=90,\qquad w_j^2+w_j\mathbin{\cdot}w=1
  \quad(1\le j\le6).
\end{equation}
\end{lemma}

\begin{proof}
Choose coordinates in which $Q(x,y)=xy$, so that
$(x,y)\mathbin{\cdot}(x',y')=(xy'+yx')/2$.  Set
\[
 h=\sqrt{90}>0,\quad R=\sqrt{1+h^2/4}>0,\quad
 b=\frac{2h}{R},\quad
 d(z)=\left(Rz-\frac h2,\frac Rz-\frac h2\right).
\]
Direct expansion gives
\[
 d(z)^2+d(z)\mathbin{\cdot}(h,h)=R^2-h^2/4=1.
\]
For a sufficiently small general nonzero complex number $\lambda$,
let $z_1,z_2,z_3$ be the roots of
\[
 X^3-bX^2+b\lambda X-\lambda=0.
\]
Vieta's formulas give
\[
 \sum_{i=1}^3z_i=b,\qquad
 \sum_{i=1}^3z_i^{-1}=b,\qquad
 2\sum_{i=1}^3d(z_i)=(h,h).
\]
Thus taking each $d(z_i)$ twice gives \eqref{curve:vector-conditions}.

For completeness, these vectors have the required distinct directions.
As $\lambda\to0$, two roots approach zero and the third approaches
$b$.  The roots are distinct for general $\lambda$: the cubic
discriminant has nonzero leading term $-4b^3\lambda$ at zero.
Near zero the projective ratio of the two entries of $d(z)$ is
\[
 \frac{z(Rz-h/2)}{R-hz/2},
\]
whose derivative at zero is $-h/(2R)\ne0$.  The two small roots
therefore give distinct nonisotropic directions.  At $z=b$ the entries
are $3h/2$ and $(R^2-h^2)/(2h)$, both nonzero, so the third limiting
direction is different from their common limit.  These properties
persist for sufficiently small general $\lambda$.  Finally apply
simultaneously to all six vectors an isometry
$(x,y)\mapsto(\nu x,\nu^{-1}y)$.  On nonisotropic projective
directions this multiplies the ratio by $\nu^2$.  Only finitely many
values of $\nu\in\CC^*$ are excluded by the prescribed finite set,
and all equations in \eqref{curve:vector-conditions} are preserved.
\end{proof}

Apply Lemma~\ref{curve:six-vectors} to the quadratic form $Q$, excluding
the branch values of the $\Phi$-map.  Each of its three directions is
then $[\operatorname{ev}_q]$ at two distinct points of $\Gamma$.
Label these points $q_1,\ldots,q_6$ so that
$[w_j]=[\operatorname{ev}_{q_j}]$, and put
\begin{equation}\label{curve:marked-divisor}
 Z=q_1+\cdots+q_6\in|3\Phi|=|M|.
\end{equation}
Choose a section $\beta$ of $M$ with divisor $Z$.
Nonisotropy implies that $F|_\Gamma$ is unramified at all six points.
Their $F$-values are not all equal, since a degree-two fiber cannot
contain six distinct points.

\subsection{Fourth-order contact on the quadric}

We have chosen the six points and their quadratic-space data on the
elliptic curve. We now construct the curve $C$ through those points.
Fourth-order contact will transfer the short evaluation schemes from
$C$ to $\Gamma$, where their section dimensions are easy to compute.

Write $\gamma$ for the section of $\OO_\Sigma(\Gamma)$ defining
$\Gamma$.  We choose a section
\[
 \kappa\in H^0(\Sigma,2\Gamma+4F),\qquad
 \kappa|_\Gamma=\beta^4,
\]
using the fixed identification in \eqref{curve:bundle-identities}.
Such sections exist: restriction to $\Gamma$ has kernel
$\OO_\Sigma(\Gamma+4F)=\OO_\Sigma(6,2)$, whose first cohomology
vanishes.  The affine space of choices is
\[
 \kappa_0+\gamma H^0(\Sigma,\Gamma+4F).
\]
A general choice defines a smooth curve $C$.  Indeed, the associated
linear system is without base points off $\Gamma$, so Bertini's
Theorem gives smoothness there.  On $\Gamma\setminus Z$ every choice
is nonzero.  At each $q_j$, global generation of
$\OO_\Sigma(6,2)$ allows the normal derivative to vary nontrivially.
Avoiding six proper affine hyperplanes makes all those derivatives
nonzero.  Their complements and the Bertini open meet in the same
irreducible affine parameter space; independence of the six evaluation
functionals is unnecessary.

More explicitly, near $q_j$ choose a local equation $y=0$ for
$\Gamma$, a coordinate $z$ on $\Gamma$ vanishing at $q_j$, and frames
in which $\kappa$ has coefficient $k(y,z)$ and $\beta$ has coefficient
$b(z)$.  The frames may be chosen compatibly with
\eqref{curve:bundle-identities}, so that
\[
 k(0,z)=b(z)^4,\qquad b'(0)\ne0,\qquad k_y(0,0)\ne0.
\]
The implicit function theorem expresses $C$ as $y=\varphi(z)$, where
\begin{equation}\label{curve:contact-expansion}
 \varphi(z)=-\frac{b'(0)^4}{k_y(0,0)}z^4+O(z^5).
\end{equation}
Thus the contact order is exactly four and
$\operatorname{div}_C(\gamma|_C)=4Z$, with the same six marked points
now regarded as points of $C$.

The bidegree of $C$ is $(8,4)$.  The sequence for its defining divisor,
together with $H^1(\Sigma,\OO_\Sigma(-8,-4))=0$, gives
$H^0(C,\OO_C)=\CC$.  Hence the smooth curve $C$ is connected and
irreducible.  Put
\begin{equation}\label{curve:line-bundles}
 H=F|_C,\qquad L=H+Z,\qquad\theta=2L,
\end{equation}
and let $\xi$ be the canonical section of $\OO_C(Z)$.
The divisor equality above specifies the isomorphism
$4\OO_C(Z)\simeq\OO_C(\Gamma)$ by $\xi^4\mapsto\gamma|_C$.
Together with adjunction, this gives
\begin{equation}\label{curve:canonical-root}
 4L=4H+4Z\simeq(\Gamma+4F)|_C\simeq K_C.
\end{equation}
We fix the last isomorphism by the following normalization.  For a
local section $s$ of $\Gamma+4F$, it is
\begin{equation}\label{curve:adjunction-normalization}
 s|_C\longmapsto
 \Res_C\left(\frac{s\gamma\Omega}{\kappa}\right),
\end{equation}
with the normal differential first.  This formula fixes the scalar in
\eqref{curve:canonical-root} for the later residue computations.

\subsection{Section counts and evaluation maps}

We give the jet calculation underlying all the special-divisor
dimensions.  The section spaces below are algebraic section spaces;
equivalently, they are holomorphic section spaces on these projective
curves, by the comparison theorem of Serre~\cite[Theorem~1]{SerreGAGA1956}.

\begin{lemma}\label{curve:adjoint-jets}
For $0\le j\le3$, set
\[
 \mathcal A_j=\Gamma+(4-j)F=(6-j,2),\qquad
 N_j=\mathcal A_j|_\Gamma=(8-j)F+\Delta.
\]
Restriction identifies $H^0(\Sigma,\mathcal A_j)$ with
$H^0(C,K_C-jH)$ and maps it onto $H^0(\Gamma,N_j)$.
For $1\le\ell\le4$, evaluation on the length-$\ell$ subschemes at
the six marked points has the same rank on $C$ as on $\Gamma$.
The ranks needed here are
\begin{equation}\label{curve:residual-table}
\begin{array}{c|c|c|c|c|c}
 j&\ell& h^0(\Gamma,N_j)&N_j-\ell Z
   &h^0(\Gamma,N_j-\ell Z)&\text{evaluation rank}\\ \hline
 0&1&16&5F&10&6\\
 1&2&14&\Phi&2&12\\
 2&2&12&\Delta&0&12\\
 3&2&10&-F+\Delta&0&10
\end{array}
\end{equation}
\end{lemma}

\begin{proof}
The restriction kernels on $\Sigma$ are respectively
\[
 \mathcal A_j-C=(-2-j,-2),\qquad
 \mathcal A_j-\Gamma=(4-j,0).
\]
Product cohomology on $\PP^1\times\PP^1$ gives $H^0=H^1=0$ for
the first kernel and $H^1=0$ for the second.  Adjunction therefore
gives the two assertions about restriction.

By \eqref{curve:contact-expansion}, the ideals of the two length-$\ell$
subschemes in the local surface are
\[
 (y-\varphi(z),z^\ell)=(y,z^\ell)\qquad(\ell\le4).
\]
Thus these are the same embedded subscheme, and restrictions of the
surface line bundle and its sections to them agree.  The two
surjectivity statements just proved imply equality of evaluation ranks.

On the elliptic curve, $\deg F=2$, $Z=3F+\Delta$, and $2Z=6F$.
These identities give every residual bundle in
\eqref{curve:residual-table}.  A positive-degree line bundle on an
elliptic curve has as many sections as its degree; a negative-degree
bundle has no section.  The degree-zero bundle $\Delta$ is nontrivial,
so both its zeroth and first cohomology vanish.  Subtracting the
residual section dimensions from $h^0(\Gamma,N_j)$ gives the stated
ranks.  In particular, the twelve-dimensional surface system for
$j=3$ has a two-dimensional kernel on restriction to $\Gamma$, and
the resulting ten-dimensional system injects into its length-twelve
evaluation space.
\end{proof}

\begin{proposition}\label{prop:curve-data}
The construction above produces a smooth connected projective curve
$C$, a reduced divisor $Z=q_1+\cdots+q_6$, line bundles $H,L,\theta$,
and a basis $u=(u_1,u_2)$ of $H^0(C,H)$ with the following properties.
\begin{enumerate}
\item The numerical and canonical data are
\[
 \begin{gathered}
 g(C)=21,\qquad \deg H=4,\qquad\deg L=10,\\
 L=H+Z,\qquad\theta=2L,\qquad K_C=4L,
 \end{gathered}
\]
where the last isomorphism is normalized by
\eqref{curve:adjunction-normalization} and $\xi^4\mapsto\gamma|_C$.
The curve $C$ has exactly fourth-order contact with the auxiliary
elliptic curve $\Gamma$ at every $q_j$.
\item The pencil $u$ defines a finite degree-four morphism
$f:C\to\PP^1$ unramified at every $q_j$, and the values $f(q_j)$ are
not all equal.  The zero divisor of $u_1$ can be chosen to be a reduced
fiber disjoint from $Z$.
\item The section dimensions are
\begin{equation}\label{curve:section-dimensions}
\begin{gathered}
 h^0(C,jH)=j+1\quad(0\le j\le3),\\
 h^0(C,L)=h^0(C,H+2Z)=2,\qquad
 h^0(C,\theta)=3,\qquad h^0(C,\OO_C(Z))=1,\\
 h^0(C,\theta+H)=6.
\end{gathered}
\end{equation}
More precisely,
\begin{equation}\label{curve:section-multiplication}
\begin{split}
 H^0(C,L)&=\xi H^0(C,H),\\
 H^0(C,H+2Z)&=\xi^2H^0(C,H),\\
 H^0(C,\theta)&=\xi^2H^0(C,2H).
\end{split}
\end{equation}
In particular $r=\xi u$ is a basis of $H^0(C,L)$ with simple common
zero divisor exactly $Z$, and every section of $\theta$ vanishes
twice along $Z$.
\item For every $k\ge2$, the evaluation map
\begin{equation}\label{curve:evaluation-surjectivity}
 H^0(C,(k+1)L)\longrightarrow H^0(Z,((k+1)L)|_Z)
\end{equation}
is surjective.
\item The auxiliary data on $\Gamma\subset\Sigma$, including the
compatible bundle identities, the logarithmic volume $\Omega$, and
the quadratic form $Q$, are as in \eqref{curve:bundle-identities} and
\eqref{curve:quadratic}; the divisor $Z$ belongs to $|M|$.
The six nonzero vectors $w_j\in H^0(\Gamma,\Phi)^*$ have the
projective evaluation directions at $q_j$, use three distinct
directions, and satisfy \eqref{curve:vector-conditions}, with
$w^2=90=9\deg L$.
\end{enumerate}
\end{proposition}

\begin{proof}
The existence and contact assertions were established above.  From
$C=(8,4)$ and $K_C=(6,2)|_C$, intersection on the quadric gives
\[
 \deg H=4,\qquad \deg K_C=6\cdot4+2\cdot8=40,
 \qquad g(C)=21.
\]
Since $\deg Z=6$, one has $\deg L=10$.  The morphism given by $H$ is
the first projection, restricted to the irreducible curve $C$, hence
is finite of degree four.  Tangency at the marked points identifies
its differential there with that of $F|_\Gamma$, which is nonzero.
Its marked values are the corresponding $F|_\Gamma$-values, already
seen not to be all equal.  Choose a regular fiber away from these
finitely many values.  A determinant-one change of the basis $u$
makes this fiber the zero divisor of $u_1$.  Such a change preserves
$u_1\,du_2-u_2\,du_1$ on both curves, and hence preserves the
quadratic form and every normalization already chosen.

For $0\le j\le3$, the sequence restricting $\OO_\Sigma(j,0)$ to
$C$ has kernel $\OO_\Sigma(j-8,-4)$, with $H^0=H^1=0$ by product
cohomology.  This proves $h^0(C,jH)=j+1$.  For any line bundle $B$
on $C$ and effective divisor $D$, Riemann--Roch and Serre duality give
\begin{equation}\label{curve:dimension-formula}
 h^0(C,B+D)-h^0(C,B)
 =\deg D-\operatorname{rank}\left(
 H^0(C,K_C-B)\longrightarrow H^0(D,(K_C-B)|_D)\right).
\end{equation}
Apply this with $B=jH$ and use the ranks in
Lemma~\ref{curve:adjoint-jets}.  Its first three rows give
\[
 h^0(C,\OO_C(Z))=1,
 \quad h^0(C,H+2Z)=2,
 \quad h^0(C,2H+2Z)=3.
\]
The inclusions of $H^0(C,H)$ into the spaces for $H+Z$ and $H+2Z$
then give $h^0(C,L)=2$.  The fourth row gives
\[
 h^0(C,3H+2Z)=h^0(C,3H)+12-10=6.
\]
This proves \eqref{curve:section-dimensions}, and equality of the
dimensions in the natural multiplication inclusions proves
\eqref{curve:section-multiplication}.  Because the pencil $u$ has no
base point and $Z$ is reduced, the assertion about the common zeros
of $r$ follows as well.

Finally, the divisor sequence for $N=(k+1)L$, followed by Serre
duality, shows that \eqref{curve:evaluation-surjectivity} is
surjective if and only if the natural inclusion
\[
 H^0(C,(3-k)L)\longrightarrow H^0(C,(3-k)L+Z)
\]
is an isomorphism.  For $k=2$ this is the inclusion from $L$ to
$L+Z=H+2Z$, between two-dimensional spaces.  For $k=3$ it is the
inclusion from $\OO_C$ to $\OO_C(Z)$, between one-dimensional
spaces.  For $k\ge4$ both spaces vanish, since the target bundle has
degree $10(3-k)+6\le-4$.  The final assertion follows from
Lemma~\ref{curve:six-vectors} and \eqref{curve:marked-divisor}.
\end{proof}

For later use, Riemann--Roch and Serre duality also give
\begin{equation}\label{curve:boundary-cohomology}
 h^1(C,L)=h^0(C,3L)=12,\qquad h^0(C,4L)=21.
\end{equation}
Indeed, $h^0(C,L)-h^1(C,L)=10+1-21=-10$ and
$K_C-L=3L$.

One additional consequence will be useful when removing a
two-dimensional obstruction.  In local frames compatible with
$L=H+Z$, with $\xi$ represented by a parameter $z$ at a marked point,
one has $r=zu$ and
$r'(0)=u(0)$.  The resulting six nonzero row vectors $r'(q_j)$ span
$\CC^2$: their projective directions are the pencil values $f(q_j)$,
which are not all equal.  Changing the frames or the local parameters
only rescales each of these rows.
